% ch08_u01 —— 原书 p.367–p.369（第 8 章全部）
\SourcePage{367}

\bdsec{8.1.}
% [Tutor Note: 原书 p.367 的 8.1 节只有编号，既无标题，底本如此，保留原样。]

\bdp{8.1.1} 选 $\mathcal{L}\in Z\operatorname{tors}_{\theta}(O)$。回忆 $\lambda_{\mathcal{L}}$ 表示 $\mathcal{GR}=G(K)/G(O)$ 上相应的局部 Pfaffian 丛（见 4.6.2）。我们要证明下述陈述，它比 5.2.14 弱，并将用于定理 5.2.14 本身的证明。

\bdp{8.1.2} \emph{命题.} 对任意 $\chi\in P_{+}({}^{L}G)$ 与 $i\in\mathbb{Z}$，$\overline{U}'$-模 $H^{i}(\mathcal{GR},I_{\chi}\lambda_{\mathcal{L}}^{-1})$ 同构于 $\mathit{Vac}'$ 的若干副本的直和。

现阶段我们并不断言这些副本的个数有限。

命题 8.1.2 是下文陈述的定理 8.1.4 与定理 8.1.6 的直接推论（第一个定理是几何的，第二个定理是表示论的）。

\bdp{8.1.3} 对 $\mathcal{GR}$ 上任一 $\mathcal{D}$-模 $M$，重整化泛包络代数 $U^{\natural}$ 作用在层 $M\lambda_{\mathcal{L}}^{-1}$ 上（见 ???）。于是 5.6.9 的典范态射 $\operatorname{Der}O\to U^{\natural}$ 给出 $\operatorname{Der}O$ 在 $M\lambda_{\mathcal{L}}^{-1}$ 上的作用。按 ???，这个作用由 $\operatorname{Der}O$ 在层 $M$ 上的作用（$\operatorname{Der}O$ 的像落在 $\mathcal{GR}$ 上向量场构成的代数中，该代数作用在 $M$ 上）与 $\operatorname{Der}O$ 在 $\lambda_{\mathcal{L}}$ 上的作用（见 4.6.7）诱导。$\operatorname{Der}O$ 在层 $I_{\chi}$ 上的作用可积分为 $\operatorname{Aut}O$ 的作用。$\operatorname{Der}O$ 在 $\lambda_{\mathcal{L}}$ 上的作用来自 $\operatorname{Aut}_{Z}O$ 在 $\lambda_{\mathcal{L}}$ 上的作用（见 4.6.7）。因此 $\operatorname{Der}O$ 在 $I_{\chi}\lambda_{\mathcal{L}}^{-1}$ 上的作用可积分为 $\operatorname{Aut}_{2}O$ 的作用。
% [Tutor Note: 原书 p.367 的 8.1.3 中两处交叉引用印作「???」，是 LaTeX 未解析的引用；底本如此，保留原样。]
于是 $L_{0}\in\operatorname{Der}O$ 在 $H^{i}(\mathcal{GR},I_{\chi}\lambda_{\mathcal{L}}^{-1})$ 上的作用可对角化，其谱含于 $\frac{1}{2}\mathbb{Z}$（事实上，依 $\operatorname{Orb}_{\chi}$ 的奇偶性，它含于 $\mathbb{Z}$ 或 $\frac{1}{2}+\mathbb{Z}$）。

\bdp{8.1.4} \emph{定理.} $L_{0}$ 在 $H^{i}(\mathcal{GR},I_{\chi}\lambda_{\mathcal{L}}^{-1})$ 上的特征值 $\ge -d(\chi)/2$，其中 $d(\chi)=\dim\operatorname{Orb}_{\chi}$。

证明将在 9.1 中给出；我们还将得到对应于 $-d(\chi)/2$ 的特征空间的如下描述。置 $F_{\chi}:=\overline{\operatorname{Orb}}_{\chi}\setminus\operatorname{Orb}_{\chi}$，$U_{\chi}:=\mathcal{GR}\setminus F_{\chi}$。$I_{\chi}$ 到 $U_{\chi}$ 的限制是（右）\SourcePage{368} $\mathcal{D}$-模 $\omega_{\operatorname{Orb}_{\chi}}$ 的正向像。它包含 $\omega_{\operatorname{Orb}_{\chi}}$ 在层论意义下的正向像，故 $H^{0}(U_{\chi},I_{\chi}\lambda_{\mathcal{L}}^{-1})\supset H^{0}(\operatorname{Orb}_{\chi},\omega_{\operatorname{Orb}_{\chi}}\otimes\lambda_{\mathcal{L},\chi}^{-1})$，其中 $\lambda_{\mathcal{L},\chi}$ 是 $\lambda_{\mathcal{L}}$ 到 $\operatorname{Orb}_{\chi}$ 的限制。于是式（241）给出嵌入

\begin{equation}\tag{392}
\mathfrak{d}_{\mathcal{L},\chi}\hookrightarrow H^{0}(U_{\chi},I_{\chi}\lambda_{\mathcal{L}}^{-1})
\end{equation}

其中 $\mathfrak{d}_{\mathcal{L},\chi}$ 是 4.6.14 中构造的 $\operatorname{Aut}^{0}_{Z}O$ 的一维表示。按 4.6.15，$L_{0}$ 在 $\mathfrak{d}_{\mathcal{L},\chi}$ 上以乘以 $-d(\chi)/2$ 的方式作用。

\bdp{8.1.5} \emph{命题.} 式（392）的像含于 $H^{0}(\mathcal{GR},I_{\chi}\lambda_{\mathcal{L}}^{-1})$。它等于 $L_{0}$ 在 $H^{0}(\mathcal{GR},I_{\chi}\lambda_{\mathcal{L}}^{-1})$ 上对应于特征值 $-d(\chi)/2$ 的特征空间。

证明包含在 9.1 中。

\emph{注.} 自然映射 $\varphi:H^{0}(\mathcal{GR},I_{\chi}\lambda_{\mathcal{L}}^{-1})\to H^{0}(U_{\chi},I_{\chi}\lambda_{\mathcal{L}}^{-1})$ 是单射，因为 $I_{\chi}$ 不可约，从而态射 $f:I_{\chi}\to R^{0}j_{*}j^{*}I_{\chi}$ 是单射，其中 $j$ 表示浸入 $U_{\chi}\hookrightarrow\mathcal{GR}$。事实上，半单性定理 5.3.3(i) 蕴含 $f$ 是同构，从而 $\varphi$ 是同构。于是命题 8.1.5 的第一个陈述，除 5.3.3 的证明中那条高度非平凡的 Lusztig 定理之外，便是显然的。

命题 8.1.2 是定理 8.1.4 与下述陈述的推论，后者将在 6.2 中证明。

\bdp{8.1.6} \emph{定理.} 设 $V$ 是离散 $U^{\natural}$-模，使得

\begin{enumerate}[leftmargin=2.2em]
\item[1)] $\mathfrak{g}\otimes O\subset U^{\natural}$ 在 $V$ 中的表示可积（即它来自 $G(O)$ 的一个表示），
\item[2)] $L_{0}\in\operatorname{Der}O\subset U^{\natural}$ 在 $V$ 上的作用可对角化，且对每个 $c\in\mathbb{C}$，其谱与 $c+\mathbb{Z}$ 的交有下界。
\end{enumerate}

则 $V$ 作为 $\overline{U}'$-模同构于 $\mathit{Vac}'$ 的若干副本的直和（即同构于 $\mathit{Vac}'\otimes W$，$W$ 为某个向量空间）。

\SourcePage{369}

\emph{注.} 设 $V$ 是离散 $U^{\natural}$-模，且作为 $\overline{U}'$-模同构于 $\mathit{Vac}'\otimes W$。更内蕴地把 $V$ 写成 $\mathit{Vac}'\otimes_{\mathfrak{z}}N$，其中 $\mathfrak{z}:=\mathfrak{z}_{\mathfrak{g}}(O)$，$N:=\operatorname{Hom}_{\overline{U}'}(\mathit{Vac}',V)=V^{\mathfrak{g}\otimes O}$。按 5.6.8，$N$ 是李代数胚 $I/I^{2}$ 上的模。$U^{\natural}$-模 $V$ 可由 $(I/I^{2})$-模 $N$ 如下重构：$V$ 是 $U^{\natural}\otimes_{\mathfrak{z}}N$ 模掉由 $u\otimes n-1\otimes an$ 生成的闭 $U^{\natural}$-子模所得的商，其中 $n\in N$，$u\in U^{\flat}_{1}$，$a\in I/I^{2}$，且 $u$ 与 $a$ 在 $U^{\flat}_{1}/U^{\flat}_{0}$ 中的像相同（见 5.6.7）。
